\section{Discussion and limitations}
\label{sec:discussion}

\subsection{Design tradeoffs}
\label{sec:tradeoffs}

The framework is organized for researchers who implement and compare
oracle-based algorithms. Python interfaces describe available access
models, checked views connect them to generators, and open RIR modules
preserve the remaining implementation choices. These mechanisms have
separate responsibilities. Structural typing is a standard Python tool;
its value here is an extensible algorithm-library interface. The runtime
boundary resolves each provider once per generation and passes the
checked view to the kernel.

RIR retains concrete signatures and assembly constants. This makes
compatibility checks and stored artifacts straightforward, while placing
width, precision, and scale sweeps in the Python generation layer.
Binding can change an implementation's private workspace and resource
requirements, but cannot silently change a public signal layout or an
alpha value already used to construct the algorithm. New implementations
can introduce new open dependencies, so successful partial linking does
not imply executable closure.

\subsection{Limits of the evidence}
\label{sec:limitations}

The reference executor, independently computed mathematical quantities,
and native backend paths answer different questions. Agreement between
backends supports a shared circuit interpretation; agreement with a
finite-precision reference supports the tested numerical implementation.
Neither observation alone establishes the accuracy of a continuum
solver or convergence over all algorithm parameters. The verification
campaign in Section~\ref{sec:numverify} records its tested instances and
revision. Current workflow reports additionally associate their results
with program and data fingerprints.

QRAM is an explicit access resource. Query and write counts do not price
memory construction, loading, latency, or fault tolerance. Logical gate
counts follow the current compiler's recipes, including its documented
QROM and constant-addition inefficiencies. Native arithmetic acceleration
changes simulation cost, not the gate-level cost assigned to the
corresponding reversible circuit. Unknown oracle implementations remain
visible in open-program reports, and their unknown workspace is not
reported as zero.

The number of algorithms in the library demonstrates breadth of
implementation, but is not by itself evidence of superior accuracy or
efficiency. Solver comparisons require matching mathematical assumptions,
target error, and success criteria. Our illustrative QFVM resource
profiles hold the problem fixed while exposing different constructions;
their low-order configurations are not an accuracy-matched ranking.
Single registers remain limited to 64 bits; general measurement and
feedback are orchestrated by host-side readout plans.

\subsection{Further development}
\label{sec:future}

The next research steps are controlled precision and parameter sweeps,
more efficient reversible arithmetic and data-loading realizations, and
additional tested approximation kernels. The open-program workflow
provides a place to compare these alternatives without discarding the
algorithm structure. Interoperability with existing fault-tolerant
algorithm libraries is also valuable: another framework's component can
be a source of an implementation when its interface and numerical
conventions are made explicit. Physical resource modeling and general
mathematical guarantees require additional methods beyond the present
logical-cost and numerical-witness facilities.

\section{Conclusion}
\label{sec:conclusion}

\pqlang{} provides a modular implementation framework for quantum
scientific-computing research. Its access interfaces, stored open
programs, partial linking, QRAM resources, and reversible numerical
computations connect paper-level oracle assumptions to concrete
implementations. Numerical witnesses and resource reports make the
consequences of those choices inspectable. The framework's contribution
is this executable and extensible research workflow, evaluated through
controlled implementation variants and deeper scientific compositions,
with the remaining numerical and physical assumptions stated explicitly.
